\scanpage{167}
% L1–8
de fonctions \unsure{semi-quadratiques|sous-quadratiques} sur $A$,
cela signifie simplement qu'elles
forment un torseur sous $L^{\otimes 2}$, par
addition de fonctions \uline{constantes}.
Si $\Lambda$ est ce torseur, on a
\[
\begin{aligned}
\Lambda\times A&\mathbin{\text{---}}\deleted{\unsure{A}}L^{\otimes 2}\\
(q,x)&\mathrel{\vdash\!\!-\!\!-}q(x)
\end{aligned}
\]
tel que
\[
(q+\alpha,x)\,\unsure{\mathrel{\vdash\!\!-\!\!-}|\mapsto}\,q(x)+\alpha\qquad\alpha\in L^{\otimes 2}
\]

% L9–13
et on prend
\[
Q_\sigma=\mathrm{Hom}_{\substack{\text{torseurs }\deleted{\text{su}}\\\text{sur }L^{\otimes 2}}}(\Lambda,L^{\otimes 2})
\]
% slanted
\margin{torseur inverse de $\Lambda$}
d'où
\[
\Lambda\simeq\mathrm{Hom}_{\text{tors sur }L^{\otimes 2}}(Q,L^{\otimes2})
\]
d'où \deleted{\unsure{accouplement|canoniquement}} application can.
\[
A\longrightarrow Q_\sigma
\]
% L13 left bubble
\margin{C'est $x\mapsto x\deleted{\unsure{\otimes}}\unsure{\cdot|\otimes}\sigma x\in\mathrm{Sym}^2(E)
\simeq\mathrm{Quad}(\underset{\substack{\Vert\\ E\otimes L^{-1}}}{\check E},k)
\simeq\mathrm{Quad}(E,L^{\otimes 2})$}

% L14–21
\deleted{Pour un choix}
Le choix d'une $q$ \inserted{$\in\Lambda$} revient à celui
d'un \unsure{isom|section} de $L^{\otimes 2}$-torseurs i.e. d'une
origine dans $Q_\sigma$ (car $Q_\sigma\simeq\Lambda$,
mais cette bijection renverse l'op.
de $L^{\otimes 2}$ en la symétrique $-$) Ceci dit,
le diviseur quadratique \inserted{$X_{\unsure{\zeta|q}}$} associé à une
$q$, correspondant à une section $q$ de $Q_\sigma$,
% archive sheet number: 166
